A popularidade de uma música pode ser examinada por medidas distintas, como posição em paradas, número de execuções e escores calculados por plataformas digitais. Essas respostas não são intercambiáveis: dependem do serviço, do período e da unidade observada. No Spotify, a popularidade é um escore de zero a cem, relacionado principalmente às execuções e à sua recência; diferentes versões de uma música podem receber valores distintos \cite{spotifydocs}. Neste estudo, interessa o escore registrado no arquivo utilizado, e não a previsão de popularidade futura ou de sucesso no lançamento.

A análise computacional da música costuma utilizar características do áudio, mas a letra oferece outra fonte de informação. Extensão, diversidade lexical e repetição podem ser descritas por medidas explícitas, enquanto representações vetoriais produzidas por modelos de linguagem sintetizam outras relações entre textos. A questão empírica é se essas representações melhoram a estimação da popularidade quando atributos acústicos e metadados já estão disponíveis.

Esse interesse tem antecedentes na recuperação de informação musical. \citeonline{dhanaraj2005} investigaram a detecção de sucessos com características acústicas e textuais. \citeonline{singhi2014} examinaram características de rima, sílabas e métrica para classificar presença em uma parada musical. Esses problemas de classificação diferem da regressão de um escore contínuo e não permitem transferir diretamente medidas de desempenho para o presente estudo. A análise de mudanças em propriedades de letras realizada por \citeonline{parada2024lyrics}, acompanhada de sua correção \cite{parada2024correction}, também motiva a investigação de descritores, sem estabelecer que eles expliquem ou determinem popularidade.

Uma dificuldade anterior ao ajuste dos modelos é a construção da base analítica. As letras não estão presentes no arquivo Spotify e sua recuperação depende de outra fonte e de regras de correspondência. Selecionar apenas músicas cuja letra foi encontrada pode modificar a distribuição de popularidade, gênero e características acústicas. Assim, a pergunta sobre ganho preditivo precisa ser acompanhada de outra: a quais faixas os resultados podem se referir?

O percurso deste trabalho parte da caracterização de 1.159.764 faixas do catálogo acessível, passa pelo sorteio estratificado de 200.000 faixas e pela recuperação de letras em um snapshot local do LRCLIB, e chega a um domínio observado de 77.535 faixas com letras aceitas. O catálogo acessível não é tratado como censo do Spotify; tampouco o subconjunto com letras é presumido representativo de todas as faixas do catálogo. Essas distinções orientam tanto a descrição dos dados quanto a interpretação dos modelos.

\subsection{Objetivos}
O objetivo geral é avaliar a associação e a contribuição preditiva incremental de representações de letras para estimar a popularidade Spotify observada, em conjunto com atributos acústicos e metadados. Especificamente, busca-se caracterizar as fontes e o efeito da disponibilidade textual sobre a composição da base; documentar a amostragem e a integração; comparar descritores lexicais e estruturais com embeddings multilingues; e avaliar os ganhos nas mesmas observações e divisões de treinamento e teste, incluindo cauda de popularidade e subgrupos linguísticos.

A comparação utiliza a mesma família de regressão, Ridge, em quatro configurações: Baseline, Tradicional, Profundo e Combinado. A denominação Profundo se refere à origem dos embeddings; não indica uma rede neural treinada neste estudo para predizer popularidade. O erro médio absoluto é a métrica principal. O interesse é a contribuição adicional das representações dentro desse procedimento, sem atribuição causal.

Além desta introdução, a Seção~\ref{sec:teoria} apresenta os fundamentos estatísticos e computacionais. A Seção~\ref{sec:metodologia} descreve as fontes, o catálogo inicial, a amostragem, a integração e a base dos modelos. A Seção~\ref{sec:resultados} discute os resultados preditivos e seus diagnósticos. A Seção~\ref{sec:conclusao} reúne as conclusões e os limites de generalização.
